% Transcription — fonds Alexandre Grothendieck, shelfmark 51, batch 1 (pages 1-20).
% Established from the facsimile archives/batches/51/batch-01.pdf (a mirror of
% https://grothendieck.umontpellier.fr/51.pdf), per the skill
% transcribe-grothendieck. The batch file carries no cover sheet: PDF page N
% is page N of the folder (the Montpellier footer prints « 1 » ... « 20 »),
% and the archivists' pencil numbers bottom-left agree (« 1 » ... « 20 »;
% the « 7 » is a hasty stroke that reads like « x »).
%
% Pass: Opus 5.5 (claude-opus-5-5), 2026-09-23 — first pass, unchecked against the pages by a human.
%
% Nineteen of the twenty pages are transcribed. Absent: page 1, a sheet of
% buff card otherwise blank, inscribed top right in his hand « "Bonne
% réduction" (Shimura, Koizumi) » — a cover, whose words become the \section
% heading (references/hand.md §5). Every other page is his mathematics; there
% are no unrelated versos in this batch. He does not paginate.
%
% The hand is his fastest: connective words are often reduced to a stroke,
% so the formulas and the nouns carry and many of the small words are read
% from the argument. The apparatus is correspondingly dense; the statements
% (Prop. 1-3, the corollaries, the two « Th. ») come through, the prose
% between them less surely.
%
% What the batch is.
%   2-5    « Proposition 1 » (V a DVR with field K, X a prescheme over V,
%          X' = X_K, F quasi-coherent on X, F' = F|X'): the quasi-coherent
%          quotients of F' correspond to the V-flat quotients of F — the
%          smallest quotient G = F/J* of F inducing G' on X' is V-flat and
%          restricts to G'; proof by π-torsion. Cor. 1 (bijection), Cor. 2
%          (compatible with ⊗_V), Cor. 3 (V-flat closed subpreschemes Y of X
%          ↔ closed subpreschemes Y' of X_K, Y = closure of Y', compatible
%          with fibre products), Cor. 4 (the closure of a closed subgroup
%          H' ⊂ G_K is the unique V-flat subgroup scheme inducing H').
%   6-7    « Prop. 2 »: G a group scheme over S locally noetherian with
%          Hilb_{G/S}; the subfunctor of closed subschemes proper and flat over
%          S that are sub-group-schemes is representable by M, with M → Hilb
%          a closed immersion; proof via a closed subscheme T = T_1 ∩ T_2 of S
%          (stable under symmetry, P = p^{-1}(H) dominating H × H).
%   8-11   « Proposition 3 »: an isogeny u': A_K → B' from the generic fibre
%          of an abelian scheme A over V extends to an abelian scheme B and
%          a homomorphism u: A → B (quotient of A by the closure of Ker u').
%          Cor. 1: good reduction is invariant under isogeny; « Théorème
%          (Shimura-Koizumi) »: a quotient of an abelian variety with good
%          reduction has good reduction — via idempotents p, q of A lifting
%          those of A', A = Ker p × Ker q. Page 11: a « Remarque » on
%          abelian subschemes B' ⊂ A_K, closure of B', B → A.
%   12-15  Counterexample in characteristic p (two finite subgroups Γ_1, Γ_2
%          of an elliptic curve B, Γ_1' ∩ Γ_2' = (e) but Γ_1 ∩ Γ_2 ≠ (e),
%          A = B/Γ_1 × B/Γ_2): the monomorphism on the generic fibre is not
%          one on the special fibre. In characteristic 0, a « Th. » (page
%          14): over S connected normal of char. 0 with Hilb_{A/S}
%          locally of finite type, every subgroup scheme B' of A_K extends
%          uniquely to a « multiabélien » subgroup B of A. A « Question »
%          (largest abelian subscheme), a remark on Hom_{S-gr}(A, B), and a
%          bracketed lemma: a rational section of a finite unramified
%          X' → X, X locally noetherian unibranch, is everywhere defined.
%   16-17  « Proposition »: A_K « se réduit bien » when it is generated by a
%          rational map from the generic fibre of a projective X over V with
%          fibres simple in codimension ≤ 1 (via a generic curve C, the
%          Jacobian Pic^0(C_K) → A_K, and Shimura-Koizumi); « Th. »:
%          Pic^0(X_K) (reduced) and Alb^0(X_K) then have good reduction.
%   18-20  « Théorème »: for f: X → S projective, S locally noetherian and
%          reduced, there is a unique abelian subpreschema A of Pic_{X/S}
%          whose underlying set is the union of the neutral components of
%          the fibres; proof via a modular problem M → S whose geometric
%          fibres are points (rigidity) and which is proper (valuative
%          criterion, Koizumi's theorem, the Chow-Murre connectedness
%          theorem). Page 20 breaks off mid-argument; page 21 (batch 2)
%          continues « Donc l'image de A est P. »
%
% Notation fixed for the batch. V the discrete valuation ring, K its field,
% k its residue field, π a uniformiser; X' = X_K, F' = F|X', and the primes
% for everything on the generic fibre as he writes them. F, G, J, N are his
% sheaves, set in italic (he writes them in a heavier upright hand). His
% « Hilb » and « Pic » are underlined twice (the representing scheme); set
% as \mathrm{Hilb}, \mathrm{Pic} without the underlining. \mathcal{M} for his
% cursive M, the scheme of the modular problem (pages 6, 13, 19-20). His
% « q. coh. », « loc. noeth. », « ss-groupe », « s-schéma », « V.A. », « ssi »
% are kept. Square brackets in the text are his.
%
% The dating: s.d. for the folder; the group « Variétés abéliennes » (45-56,
% id "45-54" in the catalogue) carries its own range 1961-1973, given as the
% weaker evidence. The quotes in \foldertitle are the inventory's own.
\documentclass[11pt,a4paper]{article}
% Le préambule commun à toutes les transcriptions du fonds Grothendieck.
%
% Il tient en une idée : **l'incertitude se compose, elle ne se résout pas.**
% Une transcription qui lisse un mot douteux a détruit la seule chose pour
% laquelle on la faisait — un lecteur doit pouvoir savoir, à chaque endroit,
% ce qui était sur le feuillet et ce qui a été ajouté. Les six macros qui
% suivent sont donc l'appareil critique, et non de la décoration.
%
% Les mêmes commandes sont reconnues par `scripts/render.mjs`, qui produit la
% vue de lecture du site. Une macro ajoutée ici doit l'être aussi là-bas,
% faute de quoi le rendu échouera bruyamment — ce qui est voulu : un rendu
% silencieusement faux est pire qu'un rendu absent.

\ProvidesPackage{grothendieck}[2026/08/08 Transcriptions du fonds Grothendieck]

\usepackage{amsmath,amssymb,amsthm}
\usepackage{mathrsfs}
\usepackage{stmaryrd}
% Les diagrammes commutatifs sont partout dans ce fonds : c'est la langue
% même de la géométrie algébrique de Grothendieck, et une transcription qui
% les rendrait en prose ne transcrirait rien.
\usepackage{tikz-cd}
% Français par défaut — c'est la langue des feuillets — mais l'anglais est
% chargé aussi : la traduction et le résumé sont des documents anglais, et la
% typographie française (espaces avant les deux-points) y serait une faute.
% Ces documents basculent par \selectlanguage{english} en tête de corps.
\usepackage[english,french]{babel}
\usepackage{fontspec}
\usepackage{geometry}
\geometry{a4paper,margin=25mm,marginparwidth=28mm}
\usepackage{marginnote}
\usepackage{xcolor}
% ulem, not soul. `soul`'s \st and \ul are built on a character-by-character
% scan that XeLaTeX's font machinery breaks: the document fails with an
% undefined control sequence rather than degrading. ulem does the same two jobs
% and is Unicode-engine safe. `normalem` stops it hijacking \emph, which the
% transcriptions use for Grothendieck's own emphasis.
\usepackage[normalem]{ulem}
\usepackage{hyperref}
\hypersetup{colorlinks=true,linkcolor=black,urlcolor=black,pdfborder={0 0 0}}

% --- Métadonnées du lot -----------------------------------------------------
% Reprises telles quelles par le script de rendu, qui les affiche en tête de
% la vue de lecture. Elles fixent ce que le fichier prétend couvrir : sans
% elles, une transcription flotte sans point d'attache dans un fonds de seize
% mille feuillets.

\newcommand{\folder}[1]{\gdef\grothFolder{#1}}
\newcommand{\batch}[1]{\gdef\grothBatch{#1}}
\newcommand{\leaves}[2]{\gdef\grothFirst{#1}\gdef\grothLast{#2}}
\newcommand{\foldertitle}[1]{\gdef\grothFoldertitle{#1}}
\gdef\grothFolder{?}\gdef\grothBatch{?}\gdef\grothFirst{?}\gdef\grothLast{?}\gdef\grothFoldertitle{}

% --- Le résumé --------------------------------------------------------------
%
% Il ouvre la lecture modernisée, sous le titre « Résumé », et il est composé
% en petit. Ce n'est pas un ornement : c'est le seul passage du document qui
% ne soit pas des mathématiques, et un lecteur doit pouvoir le reconnaître
% avant de l'avoir lu — puis le sauter s'il connaît déjà le sujet, ou s'y
% arrêter s'il ne le connaît pas. Le filet qui le suit marque la reprise du
% corps.

\newenvironment{resume}
  {\small\setlength{\parskip}{0.45em}}
  {\par\medskip\hrule\medskip}

% --- Mots-clés ---------------------------------------------------------------
%
% En anglais, en clôture du résumé de la lecture modernisée : le vocabulaire
% moderne sous lequel chercher ce que les feuillets construisent. Le manifeste
% du site (`npm run manifest`) les extrait pour étiqueter la cote entière —
% cette ligne est la source unique des tags, il n'y a pas de fichier à côté.
\newcommand{\keywords}[1]{\par\smallskip\noindent
  {\footnotesize\textsc{Keywords} --- \textit{#1}}\par}

% --- L'appareil critique ----------------------------------------------------

% Le numéro de feuillet, en marge. C'est l'ancre qui permet de régler un
% désaccord sur une formule en regardant *un* feuillet plutôt qu'une chemise
% de six cents. Le site s'en sert aussi pour tourner les pages du fac-similé
% quand on fait défiler la transcription.
\newcommand{\leaf}[1]{%
  \marginnote{\footnotesize\color{gray!70}\textsf{#1}}%
  \label{leaf:#1}\ignorespaces}

% Illisible : on ne devine pas. Un mot inventé qui se lit comme les autres est
% le pire résultat possible.
\newcommand{\ill}{\textcolor{red!60!black}{[\,\ldots\,]}}

% Lecture douteuse : le mot est donné, et signalé comme incertain.
\newcommand{\uncertain}[1]{\uline{#1}}

% Ajout de l'éditeur — une lettre manquante, un mot sous-entendu.
\newcommand{\add}[1]{\textcolor{blue!50!black}{[#1]}}

% Biffé par Grothendieck lui-même. Ce qu'il a rayé fait partie du document :
% une rature dit souvent où la pensée a changé de direction.
\newcommand{\struck}[1]{\sout{#1}}

% Note du transcripteur, sur l'état matériel du feuillet.
\newcommand{\note}[1]{\par\smallskip{\small\color{gray!60!black}\textit{[#1]}}\par\smallskip}

% Note marginale de Grothendieck, distincte du corps du texte.
\newcommand{\marginal}[1]{\par\smallskip\noindent\hspace{1em}%
  {\small\color{gray!60!black}#1}\par\smallskip}

% --- Titre ------------------------------------------------------------------

\AtBeginDocument{%
  \begin{center}
    {\large\bfseries Fonds Alexandre Grothendieck --- cote n\textsuperscript{o}~\grothFolder}\\[2pt]
    {\itshape\grothFoldertitle}\\[4pt]
    {\small feuillets \grothFirst--\grothLast\ (lot \grothBatch)}\\[2pt]
    {\footnotesize\color{gray!60!black}Transcription établie d'après le fac-similé numérisé
     par l'Université de Montpellier.\\
     Les lectures incertaines sont soulignées, les illisibles notées [\,\ldots\,],
     les ajouts entre crochets.}
  \end{center}
  \vspace{1em}}

\endinput


\folder{51}
\batch{1}
\pages{1}{20}
\dating{s.d. — le groupe « Variétés abéliennes » (45 à 56) est daté 1961-1973}
\watermark{Édition de démonstration}
\foldertitle{« Bonne réduction » (Shimura-Koizumi) : notes manuscrites (s.d.).}

\begin{document}

\section{« Bonne réduction » (Shimura, Koizumi)}
\note{titre de sa main, en haut à droite du feuillet 1, une chemise de
carton par ailleurs vierge qui ne reçoit pas de numéro de page ; les
guillemets sont les siens}

\page{2}
\marginal{Shimura -- Koizumi}
\note{encadré, en haut à droite du feuillet}

\textbf{Proposition 1.} Soient $V$ un anneau de valuation discrète
\add{interligne : de corps $K$}, $X$ un préschéma sur $V$
\add{interligne : $X' = X_K$ ouvert de $X$}, $F$ un faisceau
quasi-cohérent sur $X$, \add{interligne : $F' = F_K = F|X'$}
\struck{\uncertain{Pour que} $F$ \uncertain{soit} plat sur $V$, il f.\ et s.\
que le plus petit \uncertain{sous-schéma} \ill{} q.\ coh.\ $G$ de $F$ soit
plat sur $V$, il}
\note{quatre lignes biffées de traits horizontaux et d'une grande boucle ;
la lecture en est incertaine d'un bout à l'autre}

\struck{\ill{}} (i) Pour tout faisceau quotient quasi-cohérent
$G' = F'/J'$ \struck{$F_K$}, soit $G$ le faisceau quotient \struck{de}
$F/J^{*}$, où $J^{*}$ est le plus grand \struck{\uncertain{plus petit}}
\add{interligne : sous-Module} \struck{\ill{}} de $F$ induisant
\struck{$J'$} sur $X' = X_K$ un \struck{\ill{}} Module $\subset J'$, --- donc
$G$ est le plus petit \add{interligne : Module} quotient de $F$ induisant sur
$X'$ \struck{\uncertain{un faisceau}} un faisceau quotient de $F'$
\struck{plus grand} \uncertain{majorant} $G'$. Alors $G$ est
quasi-cohérent et $V$-plat, et $G|X' = G'$.

(ii) \struck{Soit $G = F/$ \ill{}} Soit $G = F/J$ un faisceau quotient
\uncertain{q.\ coh.}\ de \struck{$F$} $F$ qui est $V$-plat, et soit
$G' = G|X'$ le quotient correspondant de $F'$. Alors $G$ est le plus petit
Module quotient de $F$ \struck{\uncertain{majorant}} qui induise $G'$ sur
$X'$.

\textbf{Corollaire 1.} Correspondance biunivoque \uncertain{entre}
l'ensemble des faisceaux quotients q.\ coh.\ \struck{\ill{}} de $F$ qui
sont $V$-plats, et des faisceaux quotients q.\ coh.\ \add{interligne :
quotients de $F'$}.

\page{3}
\textbf{Dém.} On sait déjà \uncertain{(bof)} que $J$ est q.\ coh.\ et
$J|X' = J'$, donc $G$ est q.\ coh.\ et $G|X' = G'$. Prouvons que $G$ est
$V$-plat, i.e.\ que l'uniformisante $\pi$ de $V$ est non diviseur de $0$
dans $G$. \struck{\ill{}} Soit \ill{} \struck{\ill{}} $\mathcal{N}$
\uncertain{le noyau} \uncertain{de} $\pi$ dans $G$, alors $G/\mathcal{N}$
est un \struck{faisceau} \add{interligne : Module} quotient de $G$ qui
induit encore $G'$ sur $X'$, donc d'après la définition minimale de $G$,
$\mathcal{N} = 0$, ce qui prouve (i).

Pour (ii), on note que si $G_1$ est le plus petit Module quotient de $G$
induisant $G'$ sur $X'$, \struck{\ill{}} et on \uncertain{sait} que
$G_1 \to$ \add{interligne : \uncertain{le noyau de $F$}}
\uncertain{majoré} par $G$, \uncertain{alors} le noyau de
\add{interligne : \uncertain{libérer}} $G \to G_1$, \uncertain{est}
\add{interligne : \uncertain{inclus}} $V$-plat sur $V$ et induit $0$ sur
$X'$ ; il suffit de prouver qu'un tel Module q.\ coh.\ est Nul. Or c'est
immédiat, puisque $N \otimes_V K = 0$. Cqfd.
\marginal{$G \to G_1$}
\note{les deux mots interlinéaires de cet alinéa ne se lisent pas sûrement ;
l'argument est clair : le noyau $N$ de $G \to G_1$ est un sous-Module du
Module $V$-plat $G$, donc sans $\pi$-torsion, et nul sur $X'$}

\page{4}
\textbf{Corollaire 2.} \struck{\ill{}} \uncertain{Supposons} \ill{} qu'on
\uncertain{ait} deux $F_1$ et $F_2$ \add{interligne : q.\ coh.\ sur $X_1$,
resp.\ $X_2$} et qu'on considère les \struck{faisceaux}
\add{interligne : Modules} quotients de $F_1$, $F_2$,
$F_1 \otimes_V F_2$. Alors la correspondance du Cor.~1 \uncertain{est}
compatible avec la formation du \struck{\ill{}} produit tensoriel de
quotients des \struck{$F_1$, $F_2$, \ill{}} \add{interligne : $F_1$ et
$F_2$} ($F'_1$ et $F'_2$) identifiés à un quotient de $F_1 \otimes F_2$
($F'_1 \otimes F'_2$).

En effet, le produit tensoriel sur $V$ de faisceaux q.\ coh.\ $V$-plats est
$V$-plat, et la \struck{\ill{}} formation du produit tensoriel commute à la
tensorisation par $K$.

\textbf{Corollaire 3.} \struck{Supposons que} Les sous-préschémas
\add{interligne : fermés} $Y$ de $X$ plats sur $V$ sont en correspondance
biunivoque avec les sous-préschémas fermés $Y'$ de $X' = X \otimes_V K$,
à $Y$ correspond $Y' = Y \otimes_V K = Y \cap X'$, et à $Y'$ correspond le
sous-préschéma \uncertain{adhérence} $Y = \overline{Y'}$. Cette
correspondance est compatible avec les produits cartésiens sur $V$.

\page{5}
\textbf{Corollaire 4.} Soit $G$ un préschéma en groupes sur $V$, $H'$ un
sous-préschéma en groupes \add{interligne : fermé} de $G_K$, $H$ le
sous-préschéma adhérence de $H'$. Alors $H$ est un sous-préschéma en groupes
de $G$, plat sur $V$, \struck{dominant} induisant $H'$, et c'est le seul.

Il est évident que $H$ est stable par symétrie ; reste à prouver que
$p \colon G \times G \to G$ induit $H \times H \to H$, i.e.\ que le
sous-schéma \add{interligne : fermé} $P$ de $G \times G$ image inverse de
$H$ par $p$ \struck{\ill{}} majore $H \times H$. Or ce sous-schéma induit
sur $X'$ un sous-schéma qui majore $H' \times H'$, donc il majore
$\overline{H' \times H'}$ qui est $H \times H$ d'après le cor.~3. Cqfd.

\page{6}
\textbf{Prop.~2.} Soit $G$ un schéma en groupes sur $S$ loc.\ noeth.,
\marginal{$G$ \uncertain{q.-projectif} sur $S$ ???}
\uncertain{on} suppose que $\mathrm{Hilb}_{G/S}$
\struck{$\mathrm{Hilb}(G/S)$} existe. On s'intéresse à l'ens.\
$F(S) \subset \mathrm{Hilb}(G/S)$ formé des sous-schémas fermés de $G$,
propres et plats sur $S$, qui sont des schémas en s-groupes. Pour tout
$S'$ sur $S$, on définit de même $F(S') \subset \mathrm{Hilb}(G'/S')$,
d'où un « sous-foncteur » de $\mathrm{Hilb}_{G/S}$. Je dis qu'il est
représentable et que le préschéma $\mathcal{M}$ sur $S$ est tel que le
morphisme induit $\mathcal{M} \to \mathrm{Hilb}_{G/S}$ soit une immersion
fermée. Si $G$ est propre sur $S$ et $\mathrm{Hilb}_{G/S}$ est réunion
disjointe d'ouverts \add{interligne : $H_i$} \struck{\ill{}} de type fini
sur $S$, \struck{alors les $H_i \cap \mathcal{M} = \mathcal{M}_i$}
$\mathcal{M}$ est réunion disjointe des ouverts
$\mathcal{M}_i = \mathcal{M} \cap H_i$, dont les $\mathcal{M}_i$ sont
propres sur $S$.
\note{la marge porte, à gauche de la deuxième ligne, une note de sa main,
reproduite ci-dessus ; une flèche la rattache à « on suppose »}

N.B. Si on s'intéresse à la classification des \uncertain{s-schémas}
\add{interligne : en groupes $H$} de $G$ qui sont \uncertain{finis}
\add{interligne : \uncertain{de type} (déjà propres !)}
(resp.\ \uncertain{séparables}, resp.\ abéliens), on trouve des
sous-schémas ouverts de $\mathcal{M}$, le premier étant aussi fermé.

\page{7}
Les deux premières assertions résultent du fait suivant : Soit $H$ un
sous-schéma fermé de $G$, \add{interligne : propre et} plat sur $S$. Alors
il existe un \struck{plus petit} sous-schéma fermé $T$ de $S$ tel que pour
tout morphisme $S' \to S$, l'image inverse $H'$ \uncertain{soit} un
ss-groupe de $G'$ ssi $S' \to S$ se \uncertain{factorise} par $T$. Or, il
faut exprimer, \struck{si $H$ est} 1°) \uncertain{la symétrie},
\ill{} de $H$ et $\check{H}$ \uncertain{sont les mêmes} 2°) l'
\struck{image} image inverse de $H \times H$ est majorée par celle de
$P$, où $P = p^{-1}(H)$, $p \colon G \times G \to G$ étant la
multiplication. Les deux \struck{\ill{}} conditions en question, il est
connu, s'expriment par la \struck{\ill{}} domination \uncertain{de}
\ill{} \uncertain{par} \uncertain{des} \uncertain{sous-schémas} fermés
$T_1$, $T_2$ de $S$, on prendra donc $T = T_1 \cap T_2$.
\note{« sous-schémas » : le mot commence par ce qui ressemble à « K- »}

Pour la dernière assertion, la \uncertain{caractérisation}
\uncertain{habituelle} des propriétés \ill{} \uncertain{comme en}
Cor.~4 de prop.~1. OK.

\page{8}
\textbf{Corollaires}
\note{titre isolé en tête du feuillet, souligné ; aucun corollaire ne suit
immédiatement}

\textbf{Proposition 3.} Soit $A$ un schéma abélien sur $V$, et
$u' \colon A' = A_K \to B'$ une isogénie, où $B'$ est un schéma abélien sur
$K$. Il existe alors un schéma abélien $B$ sur $V$, \struck{\ill{}} tel
que $B \otimes_V K \simeq B'$, et un \struck{\uncertain{et cette isogénie}}
homomorphisme \struck{$A$} de schémas abéliens $u \colon A \to B$ induisant
$u' \colon A' \to B'$.

N.B. On ne se sert pas du th.\ de Weil, disant que un $V$-morphisme
\struck{$A_K \to B_K$ \ill{}} $A' \to B'$ provient d'un morphisme
$A \to B$.

\textbf{Dém.} Soit $H'$ le noyau de $u'$, c'est un sous-groupe fini de
$A'$, donc provient d'un sous-groupe $V$-plat $H = \overline{H'}$ de $A$,
évidemment fini sur $V$, étant propre \add{interligne : et
\uncertain{quasi-fini}} sur $V$ : fibre générale finie.
\marginal{\ill{}}
Comme $A$ est projectif sur $V$ (\ill{}), on considère le
\add{interligne : schéma} quotient $A/\Gamma$, qui est un schéma en groupes
sur $V$ \uncertain{car} $A$ est plat sur $V$, plat sur $V$,
(\uncertain{à fibres} \ill{}), propre sur $V$, et tel que
\note{la lettre lue $H$, $H'$ dans les deux premières lignes de la
démonstration pourrait être la même que le $\Gamma$ qui la remplace à
partir de « $A/\Gamma$ » et au feuillet 9 ; les deux glyphes se
ressemblent. La répétition « plat sur $V$, plat sur $V$ » est sur la page}

\page{9}
le noyau de $A \to B$ soit précisément $\Gamma$,
\struck{\uncertain{Prouvons maintenant} que $B$ est un schéma abélien}
\uncertain{Alors} la formation du quotient \uncertain{commute} avec le
changement de base, donc 1) $B \otimes_V K \simeq A'/\Gamma' \simeq B'$
2) $B \otimes_V k$ est isomorphe à $A \otimes k / \Gamma' \otimes k$,
\uncertain{donc} est un schéma abélien. Donc $B$ est un schéma abélien.
Cqfd.
\note{dans 2), l'indice de $\Gamma'$ est tel qu'il est écrit ; on attendrait
$\Gamma$}

\textbf{Corollaire 1.} \struck{Soit} Si $A'$ \add{interligne : et $B'$}
sont des schémas abéliens \add{interligne : isogènes} sur $K$, alors $A'$
se réduit bien par $V$ ssi $B'$ se réduit bien par $V$. S'il en est ainsi,
les \struck{\ill{}} abéliennes réduites sont isogènes.

\struck{Corollaire 2} Conséquence : \uncertain{on voit donc} \ill{}
\uncertain{qu'une} classe de variétés abéliennes mod isogénie (p.ex.\
simple) se réduit bien ou non, et \struck{\ill{}} parlons donc de la
\struck{\ill{}} classe réduite. Bien entendu, une somme de classes de
\struck{variétés} \add{interligne : schémas} abéliens sur $K$ qui se
réduisent bien, se réduit bien. En particulier, si les composantes simples
de

\page{10}
$A'$ se réduisent bien, $A'$ se réduit bien. On va prouver la
réciproque, qui \uncertain{s'exprime} :

\textbf{Théorème (Shimura -- Koizumi).} Si $A'$ sur $K$ se réduit bien,
tout schéma abélien quotient aussi.
\marginal{On \uncertain{peut} le \uncertain{prouver} directement comme
proposition~3}
\note{la note marginale est écrite en oblique, à gauche de l'énoncé, les
mots séparés par des traits}

On peut supposer \uncertain{que} $A' = A \otimes_V K$, et le
\uncertain{quotient} en question \uncertain{facteur} direct. Il y a donc des
endomorphismes $p'$, $q'$ de $A'$ tels que
\[
  p' + q' = \mathrm{id}, \qquad p'^2 = p', \qquad q'^2 = q' ;
\]
les endomorphismes correspondants $p$, $q$ de $A$ satisfont les mêmes
conditions :
\[
  p + q = \mathrm{id}, \qquad p^2 = p, \qquad q^2 = q .
\]
Donc on a un isom.
\[
  A = \operatorname{Ker} p \times \operatorname{Ker} q = A_1 \times A_2
\]
(c'est « \uncertain{fonctoriel} » !). Donc $A_1$ et $A_2$ sont propres sur
$V$, d'ailleurs le \uncertain{produit} de leurs fibres en $k$ étant
$A \otimes k$, \uncertain{qui est} \struck{\ill{}} \add{interligne :
abélien} sur $k$, il faut que $A_1 \otimes k$ et $A_2 \otimes k$
\uncertain{soient} \ill{} \uncertain{connexes} \ill{} \ill{} \ill{}
$A_1$ et $A_2$. D'où la
\marginal{\ill{}}
\note{un long trait vertical dans la marge de gauche, en bas du feuillet,
porte quelques lettres illisibles}

\page{11}
conclusion.

\textbf{Remarque.} Soit $A$ un schéma abélien sur $V$, $B'$ un
sous-schéma abélien de $A' = A_K$. Il résulte de ce qui précède que $B'$
est de la forme $B_K$, où $B$ est un schéma abélien sur $V$, et où
l'immersion $B' \to A'$ est induite par un morphisme $B \to A$ (ce dernier
point résulte aussi du th.\ de Weil). \struck{Mais} Soit $H$ le
sous-schéma fermé de $A$ image de $B$, on voit \struck{facilement}
\add{interligne : tout de suite} que $H = \overline{B'}$ (car $H$ induit
$B'$, et donc $H$ majore $\overline{B'}$, or $B$ l'image inverse de
$\overline{B'}$ est un sous-schéma fermé de $B$ dominant $B \otimes K$, on
identifie : $B$, donc $\overline{B'}$ majore $H$). C'est donc un
sous-schéma en groupes fermé de $A$, \add{interligne : plat sur $V$}
(cf.\ prop.~1, cor.~4). C'est un schéma abélien, \uncertain{et}
$B \to B'$ est un isom.\ (grâce au th.\ de Weil p.ex.), i.e.\ $B \to A$ est
un \uncertain{monomorphisme}, i.e.\ son noyau est \struck{\ill{}}
\uncertain{trivial}.
\note{« $B \to B'$ » : sans barre sur la page ; l'argument demande
$B \to \overline{B'}$}

D'après Shimura -- Koizumi, il \uncertain{n'en est}

\page{12}
pas toujours ainsi, comme on \uncertain{voit} ainsi. \uncertain{Soit}
\struck{\ill{}} \uncertain{car.}\ $k = p > 0$, et \uncertain{prenons}
deux ss-schémas finis $\Gamma_1$, $\Gamma_2$ du schéma abélien
\struck{$A$} $B$, tels que
$\Gamma'_1 \cap \Gamma'_2 = (e)$, $\Gamma_1 \cap \Gamma_2 \neq (e)$. On en
construit facilement, p.ex.\ en partant de $B$ courbe elliptique telle que
$B \otimes k$ n'ait pas de \uncertain{points} d'ordre $p$ (donc
$B \otimes k$ a un seul sous-groupe fini d'ordre $p$) et $B \otimes K$
ayant deux sous-groupes distincts d'ordre $p$ (c'est faisable
\struck{\ill{}} en \uncertain{inégales} caractéristiques $0$, $p$,
\struck{\ill{}} \uncertain{quitte à} \struck{\ill{}} \ill{} une
extension finie de $K$).
\note{trois lignes lourdement biffées, traversées d'une ligne ondulée,
entre « caractéristiques » et « une extension finie de $K$ » ; on y lit
« $B \otimes K$ », rien de suivi}

\uncertain{On pose} $A = B/\Gamma_1 \times B/\Gamma_2$, et on y envoie $B$
diagonalement. OK.

\marginal{L'\uncertain{exemple} \uncertain{marche} \uncertain{encore}
quand pour $B$ un schéma \struck{\ill{}} abélien
\add{interligne : \ill{}} quel.\ de dim $\geq 2$, \struck{\ill{}}
\uncertain{pourvu} \uncertain{que} la puissance $p$-ième dans
$\mathrm{Lie}(B)$ soit nulle, \uncertain{de sorte} \uncertain{qu'on}
\ill{} \ill{} \ill{} $B$ \ill{} $\mathrm{Lie}(B)$ \struck{\ill{}}
\ill{} \ill{}, \uncertain{définissant} \ill{} \ill{} \ill{}}
\note{note marginale écrite en oblique dans la marge inférieure gauche,
séparée du corps par une grande courbe qui aboutit au $B$ de « on y envoie
$B$ diagonalement »}

Noter cependant que en car.~$0$ \add{interligne : \ill{}}, \uncertain{tout}
schéma en groupes de type fini sur un corps est \uncertain{séparable},
\uncertain{donc},

\page{13}
il est immédiat de prouver que \struck{\ill{}} \uncertain{tout}
sous-schéma \struck{$A$} en groupes \add{interligne : fermé} $B'$ de
$A' = A_K$ est la restriction d'un \struck{\ill{}} ss-schéma abélien $B$
de $A$. Il s'ensuit, si on suppose que $A$ est un schéma abélien
\add{interligne : \uncertain{projectif}} sur un schéma loc.\ noeth.\ $S$ de
car.~$0$, \struck{tel} tel que $\mathrm{Hilb}_{A/S}$
\struck{$\mathrm{Hilb}(A/S)$} existe \add{interligne : et soit loc.\ de
type fini} (p.ex.\ si $A$ projectif sur $S$) que le schéma $\mathcal{M}$
sur $S$ de prop.~2 \struck{est tel que} qui classifie les sous-schémas
abéliens \struck{(\ill{})} est \struck{tel que ses \ill{} (\ill{}) \ill{}}
loc.\ propre sur $S$. D'ailleurs, \ill{} \uncertain{en caractéristique}
\ill{}, le $S$-schéma $\mathcal{M}_1$ \add{interligne : en groupes finis}
qui classifie les ss-schémas (\uncertain{séparables} en
\add{interligne : loc.\ non ramifié}) \add{de} $A$ est
\struck{\ill{}} sur $S$. \uncertain{Donc}

En caractéristique~$0$ en combinant les deux propriétés, on voit ceci :
\note{la phrase commencée par « Donc » est laissée en suspens ; « En
caractéristique $0$ » la reprend à la ligne, et le feuillet s'arrête sur
les deux-points}

\page{14}
\textbf{Th.} \struck{Proposition} Soit $A$ un schéma abélien sur le
\struck{\ill{}} préschéma connexe et \uncertain{normal}\,? $S$ de car.~$0$,
de corps $K$. \struck{\uncertain{Supposons} \ill{} schéma en groupes}
\uncertain{tel} que $\mathrm{Hilb}_{A/S}$ existe et soit loc.\ de type
fini sur $S$ [p.ex.\ que $A$ soit projectif sur $S$]. Alors pour tout
sous-schéma en groupes $B'$ de $A' = A \otimes K$, il existe un unique
sous-schéma en groupes \uncertain{multiabélien} $B$ de $A$, tel que
$B' = B \otimes K$. Si $B'$ est abélien, $B$ est abélien.
\marginal{N.B. Il suffit que $A$ soit propre plat sur $S$ \ill{} \ill{}
\uncertain{la composante} \uncertain{neutre} \uncertain{de} $A$
\uncertain{multiabélienne}) \ill{} \uncertain{et} $S$ \ill{}
\uncertain{suffit}}
\note{note marginale en oblique, à gauche de l'énoncé ; « normal » est
souligné et suivi d'un point d'interrogation. « Multiabélien » revient
au feuillet 15, « schémas multi-abéliens »}

Pour \uncertain{un} \uncertain{emploi} de l'hypothèse : $S$ normal,
\uncertain{Ce n'est} \uncertain{pas} \uncertain{bien} \ill{} \uncertain{ :
$S$ réduit} ? \struck{\ill{}} probable, mais \struck{\ill{}} \ill{}
\uncertain{si} $S$ est \uncertain{quelconque}. \struck{\ill{}}
\uncertain{pour} \uncertain{que} $\mathcal{M}$ \uncertain{soit}
\uncertain{non} \uncertain{ramifié} sur $S$.
\note{alinéa très raturé : deux lignes biffées, puis un bloc encadré et
biffé d'une ligne ondulée, où se lit « unibranche » en interligne ; un
crochet et un trait dans la marge de gauche}

\textbf{Question} ? Est-il possible en car.~$0$ de définir, si $A$ est un
\struck{\ill{}} schéma en groupes propre sur $S$, un plus grand
sous-schéma abélien [\uncertain{qui} \uncertain{coïncide} \ill{} \ill{}
\uncertain{de} $A$ \uncertain{si} \ill{}]. \uncertain{Ceci si} $S$ est
réduit ?

\page{15}
\textbf{Remarque.} \struck{\ill{}} \uncertain{Pour}
$\mathrm{Hom}_{S\text{-}\mathrm{gr}}(A,B)$ [\uncertain{où} $A$ et $B$
\uncertain{sont} \uncertain{multi}abéliens sur $S$] la situation est
analogue. \uncertain{Sans} hypothèse de \struck{\ill{}} caractéristique :
ce schéma sur $S$ [\uncertain{sous réserve} d'existence, et d'être
\struck{de} \add{interligne : loc.\ de} type fini] est loc.\ ? non ramifié
et loc.\ propre sur $S$. D'où un théorème analogue au précédent sur
l'existence des $\mathrm{Hom}$ de schémas multi-abéliens.

[\textbf{Lemme.} \struck{Soit} Soit $f \colon X' \to X$ un morphisme
\uncertain{fini} et non ramifié, avec $X$ \struck{normal} loc.\ noeth.\
unibranche. Alors toute section rationnelle de $X'$ sur $X$ est partout
définie.

En effet, \ill{} on se ramène au cas où \struck{$f$} $X$ est intègre
\add{interligne : affine}, $X'$ intègre, $f \colon X' \to X$ dominant
affine non ramifié et birationnel, donc défini par une
\uncertain{algèbre} \struck{\ill{}} $A'$ sur $A$, contenue dans la clôture
intégrale de $A$ dans $K$. \struck{\ill{}} Alors les fibres
\uncertain{de $f$} \ill{} \uncertain{ont un seul} pt géométrique.
\uncertain{Étant} séparable, l'\uncertain{une} \ill{} \ill{}, et on finit :
$f$ est un isom.]
\note{les crochets ouvrant et fermant sont les siens et embrassent tout le
lemme}

\page{16}
\textbf{Proposition.} Soient $V$ un anneau de valuation discrète, de corps
$K$, $A_K$ une V.A.\ définie sur $K$, $X$ un préschéma projectif sur $V$,
et $R_1$ sur $V$ [\struck{fibres} plat, fibres simples en codim $\leq 1$],
$u \colon X_K \to A_K$ une application rationnelle qui « engendre » $A_K$.
Alors $A_K$ « se réduit bien » sur $V$, i.e.\ il existe un $A$ abélien sur
$V$ tel que $A_K$ soit isomorphe à $A \otimes_V K$.

\struck{On sait que} Soit $C$ une « courbe générique » dans
\uncertain{$X$}, \uncertain{successivement} simple sur $V$. (On sait que
cela existe…). \struck{Considérons} \struck{L'\uncertain{image}} Je
\uncertain{dis} \uncertain{qu'on} \uncertain{sait} que alors
$C_K \to A_K$ est un \uncertain{morphisme}, et qui engendre $A_K$
[\uncertain{quitte à} \ill{}, \uncertain{sur} un corps alg.\
\struck{\ill{}} \uncertain{clos}]. Considérons l'hom.\ correspondant
$\mathrm{Pic}^0(A_K) \to \mathrm{Pic}^0(C_K)$, qui
\struck{est injectif \ill{} \uncertain{fait} \ill{} \uncertain{un}
\uncertain{isomorphisme} \ill{} $\mathrm{Pic}^0(A_K)$ \ill{} se réduit
bien, \ill{}}
$\mathrm{Pic}^0(C_K) \to A_K$ qui est surjectif [\uncertain{question} \ill{}
\uncertain{sur} un corps]. Comme $\mathrm{Pic}^0(C_K)$ se réduit bien grâce
à $\mathrm{Pic}^0(C)$, il s'ensuit \uncertain{par} le th.\ de
Koizumi -- Shimura \uncertain{que} $\mathrm{Pic}\,A$ se réduit bien.
\note{« dans $X$ » : la lettre ressemble à un $V$ ; « successivement » est
une conjecture sur un mot réduit à un trait. Le bloc biffé est encadré et
couvert de boucles ; la suite, $\mathrm{Pic}^0(C_K) \to A_K$, est écrite à
droite du bloc. Le passage de $\mathrm{Pic}\,A$ à $A$ n'est pas écrit}

\page{17}
\textbf{Th.} Soient \struck{$A$} $V$ un anneau de valuation discrète, $X$
un préschéma projectif sur $V$, à fibres \struck{simples séparables}
simples en codim $\leq 1$, \add{interligne : $X_K$ géom.\
\uncertain{irréd.}} Alors $\mathrm{Pic}^0(X_K)^{\uncertain{\mathrm{red}}}$
et $\mathrm{Alb}^0(X_K)$ se réduisent bien.
\note{l'exposant de $\mathrm{Pic}^0(X_K)$ est surchargé et peu lisible}

Rappelons nous que s'il existe dans $\mathrm{Pic}^0(X_K)$ un sous-schéma
abélien $\mathrm{Pic}^0(X_K)'$, \uncertain{ayant} \uncertain{même}
ens.\ sous-jacent, et que $\mathrm{Alb}^0(X_K)$ est défini comme le dual de
ce dernier.

\struck{\ill{} $\mathrm{Alb}^0(X_K)$ \ill{} \ill{} \ill{} dual \ill{}}
\note{une ligne biffée et encadrée, couverte d'un gribouillis}
\marginal{$X_K \; x_1$}
Faisons une extension \uncertain{non ramifiée} sur $V$, on peut supposer
que $X$ a une section \uncertain{sur} $\mathrm{Spec}(V)$. Cette section
nous donne un \struck{\ill{}} morphisme naturel
$X_K \to \mathrm{Alb}^0(X_K)$, et il est connu que \struck{\ill{}} ce
morphisme engendre $\mathrm{Alb}^0(X_K)$ \struck{\ill{}} (c'est
\uncertain{formel}). Donc par la proposition précédente,
$\mathrm{Alb}^0(X_K)$ se réduit bien. Mais comme $\mathrm{Pic}^0(X_K)'$
\uncertain{est} \uncertain{isogène}, il s'ensuit que ce dernier
\struck{\ill{}} se réduit bien également.

\page{18}
\marginal{Ceci \uncertain{lorsque}, \uncertain{en fait}, on
\uncertain{prouve} \uncertain{d'abord} \ill{} que
$\mathrm{Pic}^{0}_{X/S}$ est ouvert, donc en l'\uncertain{occurrence}
\uncertain{simple} sur $S$.}
\note{note encadrée en tête du feuillet, reliée par une flèche au mot
« Théorème »}

\textbf{Théorème.} Soient $f \colon X \to S$ un morphisme projectif,
\struck{plat, \add{interligne : et \ill{} ?} à fibres géométriques
\uncertain{intègres} \ill{} et simples en codim~$1$}, le but étant loc.\
noethérien et réduit. Alors il existe un sous-préschéma abélien unique $A$
de $\mathrm{Pic}_{X/S}$, dont l'ensemble sous-jacent soit la réunion des
composantes connexes \add{\uncertain{neutres}} des fibres de
$\mathrm{Pic}_{X/S}$.
\marginal{Prouvé seulement en car.~$0$ pour l'instant…}
\marginal{est-ce nécessaire ? Non, voir \uncertain{la}
\uncertain{rigidité} des Modules formels des V.A.\ $*$}
\note{la seconde note marginale est entourée et reliée par une flèche au mot
« réduite », souligné dans l'énoncé ; « neutres » manque à la lettre et
est suggéré par la suite (« ses fibres sont propres »)}

Soit en effet $P$ ladite réunion. Comme ses fibres sont propres sur $k$
[\uncertain{vrai} …], il s'ensuit par un théorème connu
[\uncertain{renvoyant} aux critères de propreté de Ch.~II] que $P$ est
propre sur $S$ [question : en est-il de même de
$\mathrm{Pic}^{\tau}_{X/S}$ ? De façon générale, est-il vrai que
$\mathrm{Pic}_{X/S}$ est \uncertain{réunion} \uncertain{disjointe}
\uncertain{de} \uncertain{parties} \uncertain{ouvertes} propres sur $S$ ?].
\marginal{? Non \ill{} \uncertain{seulement} \uncertain{formel}…}
Nous cherchons un sous-schéma fermé de $P$, qui soit plat sur $S$ et qui
soit un sous-préschéma en groupes de $\mathrm{Pic}_{X/S}$
\add{interligne : et abélien sur $S$, ensemblistement égal à $P$}
(\struck{\ill{}} \ill{}), pour $S'$ \uncertain{sur} $S$
\uncertain{variable},
\[
  F(S') = \text{l'ens.\ des sous-schémas fermés de } P_{S'} ,
\]
qui sont des sous-schémas en groupes de $G$, [abéliens sur $S$,
ensemblistement égaux à $P_{S'}$]. Je dis que le foncteur $F(S')$
\note{« Ch.~II » : le chiffre romain est encadré ; « $G$ » est tel qu'il
est écrit, pour $\mathrm{Pic}_{X/S'}$}

\page{19}
est représentable. En effet, laissant d'abord tomber la dernière
condition, on est ramené à un sous-schéma \struck{fermé} de
$\mathrm{Hilb}_{P/S}$ [\uncertain{formé} des sous-schémas
\add{interligne : plats} de $P$ qui sont des sous-groupes de
$\mathrm{Pic}$, et \uncertain{ouvert} dans ce dernier par la condition
abélienne], montrons que la condition ensembliste est une condition
\uncertain{ouverte} : elle est en effet invariante par changement de base,
et si elle est vérifiée par un $A \subset P$ en un \uncertain{pt}, elle
l'est par raison de dimensions, dans un voisinage. De plus, le morphisme
$\mathcal{M} \to S$ de ce problème modulaire a les propriétés suivantes :

(i) \struck{Les} \uncertain{Ses} fibres géométriques sont isomorphes
\uncertain{à} \struck{\ill{}} $\mathrm{Spec}\, k(x)$.
\uncertain{Évident}, on \uncertain{se ramène} sur un corps alg.\ clos, et
utilisant le \uncertain{th.\ de rigidité}.
\note{« th.\ de rigidité » est écrit en interligne et souligné deux fois}

(ii) Le morphisme $\mathcal{M} \to S$ est propre. \uncertain{Il est} en
effet de type fini, séparé, \uncertain{reste} à vérifier le critère
valuatif, savoir ceci : dans le cas où $S$ est le spectre d'un

\page{20}
anneau de valuation discrète, et \uncertain{où} on se donne une section
\add{interligne : rationnelle} de $\mathcal{M}$ sur $S$ i.e.\ \uncertain{on}
\uncertain{se donne} l'unique $A_K$ abélien \uncertain{sous-groupe} de
$\mathrm{Pic}^0(X_K)$\,\ill{} \uncertain{prouver} que la section est partout
définie. Or par le théorème de Koizumi $A_K$ « se réduit bien », i.e.\
provient d'un schéma abélien $A$ sur $V$. On va essayer de prolonger
\struck{\ill{}} en un morphisme $A \to \mathrm{Pic}_{X/S}$, i.e.\
$A \to P$. C'est possible, vu le fait que les fibres de $P$ sur $S$
satisfont les hypothèses du théorème de connexion
\add{interligne : \uncertain{linéaire}} de Murre -- Chow. [On n'a pas
besoin du théorème si $X$ simple sur $S$].
\note{« rationnelle » est entouré ; après $\mathrm{Pic}^0(X_K)$, un signe
qui ressemble à un astérisque suivi de « + »}

Reste la question si $u \colon A \to P$ est une immersion fermée, d'image
$P$. L'image \struck{\ill{}} est évidemment
\add{interligne : \uncertain{une fibre}} \uncertain{de} dim.\ égale à
$\dim A_K$ \add{interligne : $= \dim \mathrm{Pic}(X_K)$} \struck{et}
[\uncertain{continuité} de la dimension des fibres du noyau] donc de dim
égale à $\dim \mathrm{Pic}(X_K)$ [théorème de continuité de la dimension
des fibres des Picards].
\marginal{Attention, cela est-il vrai si $f \colon X \to S$ n'est pas
simple ?!}
\note{dans le dernier $\mathrm{Pic}(X_K)$, on attendrait la fibre spéciale
$X_k$ ; l'indice est écrit comme les précédents. La démonstration continue
au feuillet 21 (batch 2) : « Donc l'image de $A$ est $P$. »}

\end{document}
